\chapter{Introduction}

Les biographies Wikipedia contiennent des informations riches mais heterogenes:
redondances, variations de style, granularite temporelle irreguliere et contenu
multilingue. Generer une narrative analogue coherente a partir de telles sources
necessite plus qu'une generation libre par LLM; il faut une chaine structuree
qui extrait, normalise, compare et valide.

L'objectif de ce projet est de construire un pipeline logiciel capable de:
\begin{itemize}[leftmargin=2em]
  \item extraire des timelines factuelles depuis du texte brut,
  \item produire une representation anglaise des timelines FR
        (\texttt{fr -> fr\_en}) pour un apprentissage cross-lingue,
  \item construire des paires supervisees \texttt{en <-> fr\_en},
  \item entrainer un modele de similarite narrative,
  \item generer des candidats analogues et les filtrer automatiquement via un seuil.
\end{itemize}

Une branche alternative de fusion semantique FR/EN puis finalisation
\texttt{final\_en} a aussi ete developpee, mais les resultats quantitatifs
retenus dans ce rapport proviennent du pipeline \texttt{en/fr\_en}.

La contribution principale est un systeme operationnel de bout en bout,
accompagne d'une evaluation experimentale de modeles de similarite.
Le rapport suit la progression du pipeline depuis la donnee brute jusqu'a
l'interface de demonstration.
